% ===================================================================
%  6. TAULA — Etxea barrutik. — Altzariak. — Janaria. — Argiztapena.
%  Traduction 2026 d'apres texto/io, controlee sur texto/fr, texto/en
%  et texto/es. Consigne : voir texto/eu/10-taula-01.tex.
%
%  Deux notes, toutes deux marquees « (*) », et deux appels : celui de
%  « babo » au bloc c1-12-1, celui de « lambrequino » au bloc c2-01-1.
%
%  LE BLOC c1-06-1 GARDE SON ECART DECLARE : l'ido y porte « (6) » la
%  ou le francais porte « (16) », et c'est le francais qui a raison.
%
%  LE PLUS GROS LEXIQUE DE LA COLONNE — cinq pieces, deux cent
%  cinquante objets — et le plus riche en composes basques : la
%  passoire est un « iragazkia », ce qui filtre ; la rape un
%  « birringailua », ce qui broie ; le tire-bouchon un
%  « kortxo-kentzekoa », ce qui ote le bouchon. La langue nomme
%  l'ustensile par sa fonction, sans exception.
% ===================================================================

%%K t06-apar-1 apar
\VUsaut{6.0mm}
\VUtitre{15.0pt}{60}{6. TAULA}
\VUsaut{1.4mm}
\VUtitre{11.4pt}{0}{Etxea barrutik, Altzariak, Janaria,}
\VUsaut{0.4mm}
\VUtitre{11.4pt}{0}{Argiztapena.}
\VUsaut{2.2mm}

%%K t06-apar-2 apar
Etxea prest dago. Joanesen osaba bere etxebizitza berrira aldatu da, zeinaren
banaketa ezagutzen baitugu jada. Ikus dezagun orain nola hornitu duen. Lehenik
hauek bisitatuko ditugu:

%%K t06-tit-1 sub
\VUpk{10.2pt}{40}{Logela.}
\VUsaut{3.0mm}

%%K t06-c1-01-1 p
1. --- Une txarrean iritsi gara, neskameak \VUgras{gela} garbitzen ari baita.

%%K t06-c1-02-1 p
2. --- \VUgras{Garbigailuaren} \textsuperscript{(1)} gainean hor-hemenka gauza
desberdinak daude, zeinekin garbitu, orraztu eta oro har txukuntzen baita.
\VUgras{Ontzia} \textsuperscript{(2)} eta \VUgras{pitxerra}
\textsuperscript{(3)}, \VUgras{ile-eskuila} \textsuperscript{(4)},
\VUgras{orrazia} \textsuperscript{(5)}, \VUgras{xaboia} \textsuperscript{(6)} eta
\VUgras{belakia} \textsuperscript{(7)} dira, eta azkenik \VUgras{flaskoa}
\textsuperscript{(8)}, ahoa garbitzeko urarekin.

%%K t06-c1-03-1 p
3. --- Zoruan, garbigailuaren aurrean, ur zikinerako \VUgras{ontzi} bat
\textsuperscript{(9)} dago.

%%K t06-c1-04-1 p
4. --- \VUgras{Atalondoan} \textsuperscript{(10)} \VUgras{esekitoki} batzuk
\textsuperscript{(13)} eta bi \VUgras{aterki} \textsuperscript{(11)} ikusten dira
\VUgras{euskarrian} \textsuperscript{(12)}.

%%K t06-c1-05-1 p
5. --- Atearen beste aldean \VUgras{ispilu-armairua} \textsuperscript{(14)} dago,
zeinetan \VUgras{oihalak} \textsuperscript{(15)} gordetzen baitira.

%%K t06-c1-06-1 p
6. --- Balia gaitezen \VUgras{neskamea} \textsuperscript{(16)} \VUgras{ohea}
\textsuperscript{(17)} egiten ari denez, eta azter ditzagun haren zatiak.
\VUgras{Ohatzeak} \textsuperscript{(20)} \VUgras{malguki-sarea}
\textsuperscript{(21)} sendoa darama, zeinaren gainean Justinak berehala jarriko
baitu artilezko \VUgras{koltxoia} \textsuperscript{(22)}. Gero aulkietatik hartuko
ditu bi \VUgras{maindire} \textsuperscript{(23)} eta \VUgras{estalkia}
\textsuperscript{(24)}. Ondoren buru aldean \VUgras{luzetarako burkoa}
\textsuperscript{(25)} eta \VUgras{burkoa} \textsuperscript{(26)} jarriko ditu,
zeinak \VUgras{lumatzarekin} \textsuperscript{(27)} batera baitaude
\VUgras{ohe-ondoko tapizaren} \textsuperscript{(32)} gainean. Erratzaz
\VUgras{gortinak} \textsuperscript{(19)} eta \VUgras{oheburuko oihala}
\textsuperscript{(18)} astinduko ditu, eta lanaren zati bat eginda egongo da.

%%K t06-c1-07-1 p
7. --- Gero \VUgras{ohe-ondoko mahaitxoa} \textsuperscript{(28)} pixka bat
txukundu beharko du, \VUgras{iratzargailua} \textsuperscript{(29)} eman,
\VUgras{gau-argia} \textsuperscript{(30)} prestatu eta \VUgras{kandela}
\textsuperscript{(31)} eraman, kandelargia garbitzeko.

%%K t06-tit-2 sub
\VUsaut{5.0mm}
\VUpk{10.2pt}{40}{Josteko saskia eta tximiniako tresneria.}
\VUsaut{3.0mm}

%%K t06-c1-08-1 p
8. --- \VUgras{Josteko saskiak} \textsuperscript{(34)}, \VUgras{taburetearen}
\textsuperscript{(33)} ondoan, hark ere garbiketa handia behar du.
\VUgras{Brodatua} \textsuperscript{(35)} tolestu beharko da, \VUgras{lan-mahaitxoan}
\textsuperscript{(38)} gorde \VUgras{titarea}, \VUgras{hariak}
\textsuperscript{(36)}, \VUgras{orratzak} \textsuperscript{(37)} eta
\VUgras{guraizeak} \textsuperscript{(39)}, eta mahaitxoaren estalkia
\VUgras{giltzaz} \textsuperscript{(40)} itxi.

%%K t06-c1-09-1 p
9. --- Gelaren erdiko \VUgras{mahai biribilean} \textsuperscript{(41)} Justinak
\VUgras{karrafaren} \textsuperscript{(42)} ura aldatuko du. \VUgras{Azukrea}
\VUgras{azukre-ontzian} \textsuperscript{(43)} jarriko du, \VUgras{azukre-matxarda}
\textsuperscript{(44)} garbituko du eta \VUgras{arropa-eskuila}
\textsuperscript{(46)} eramango.

%%K t06-c1-10-1 p
10. --- Gelaren eskuinaldeak lan gutxiago emango dio, \VUgras{sofa}
\textsuperscript{(47)} eta haren \VUgras{kuxina} \textsuperscript{(48)} beren
lekuan baitaude, eta hotzik egiten ez duenez, \VUgras{tximiniaren}
\textsuperscript{(49)} tresneria ere ez baita ukitu. \VUgras{Palatxoa}
\textsuperscript{(50)}, \VUgras{matxarda} \textsuperscript{(51)},
\VUgras{burdinsuak} \textsuperscript{(55)} eta \VUgras{hauts-sareta}
\textsuperscript{(56)} garbi daude; \VUgras{tximinia-pantaila}
\textsuperscript{(54)} ez da mugitu, eta \VUgras{egur-kutxa}
\textsuperscript{(52)} \VUgras{egurrez} \textsuperscript{(53)} beteta dago.

%%K t06-noto-1 noto
\VUnotes{143.2mm}{%
(*) Babo = haur txikia; ingel. \textit{baby}, fran. \textit{bébé}, ital.
\textit{bambino}.%
}

%%K t06-noto-2 noto
\VUnotes{143.2mm}{%
(*) Lambrequino = fran. \textit{lambrequin}, esp. \textit{lambrequines}, port.
\textit{lambrequins}, eta alemanez eta ingelesez ere \textit{lambrequins} ---
hau da, alemanez, ingelesez, frantsesez, portugesez eta gaztelaniaz.%
}

%%K t06-c1-11-1 p
11. --- Oraindik \VUgras{tximiniako erlojua} \textsuperscript{(57)},
\VUgras{kandelabroak} \textsuperscript{(58)} eta \VUgras{gauzatxoen erretiluak}
\textsuperscript{(59)} astindu beharko ditu, eta gero \VUgras{ispilua}
\textsuperscript{(60)} igurtzi.

%%K t06-c1-12-1 p
12. --- Haurtxoaren (babo (*)) \VUgras{sehaskari} \textsuperscript{(61)}
dagokionez, prest dago jada.

%%K t06-tit-3 sub
\VUsaut{5.0mm}
\VUpk{10.2pt}{40}{Egongela.}
\VUsaut{3.0mm}

%%K t06-c2-01-1 p
1. --- Hona hemen egongela. \VUgras{Gela} oso ederra da. Eskuineko izkinako
tximinia \VUgras{estatuatxo} dotore batez \textsuperscript{(1)} eta bi
\VUgras{ontzi} ederrez \textsuperscript{(3)} apainduta dago, eta belusezko
\VUgras{lanbrekinez} \textsuperscript{(2)} (*) estalita.

%%K t06-c2-02-1 p
2. --- Sutegiko txinpartek tapiza erre ez dezaten, tximiniaren aurrean kobrezko
\VUgras{pantaila} bat \textsuperscript{(4)} jarri dute.

%%K t06-c2-03-1 p
3. --- Ondoan andreen \VUgras{idazmahai} dotore bat \textsuperscript{(5)} dago
zabalik. Hemen idazten ditu \VUgras{etxekoandreak} \textsuperscript{(23)} bere
gutunak.

%%K t06-c2-04-1 p
4. --- Leihoa \VUgras{tulezko errezelez} \textsuperscript{(6)} apainduta dago,
\VUgras{lokarriz} \textsuperscript{(8)} bilduak \VUgras{kakoen}
\textsuperscript{(7)} gainean, eta oihalezko gortina astunez, goian
\VUgras{lanbrekinarekin} \textsuperscript{(9)}.

%%K t06-c2-05-1 p
5. --- Leihoaren hobian \VUgras{lore-ontziak} \textsuperscript{(11)}
\VUgras{landare} berde bat \textsuperscript{(10)} eusten du, palmondo bikain bat,
zeinaren hostoak ia \VUgras{petrolio-lanpararaino} \textsuperscript{(12)} iristen
baitira.

%%K t06-c2-06-1 p
6. --- \VUgras{Pantaila} \textsuperscript{(13)} Mariak egin zuen,
\VUgras{neskatxa} xarmangarriak \textsuperscript{(14)}, \VUgras{pianoaren}
\textsuperscript{(15)} aurrean dagoenak. \VUgras{Taburetean}
\textsuperscript{(16)} oso tente eserita, pieza bat ikasten ari da. Oso trebea eta
oso txukuna da, eta hori erakusten du bere \VUgras{partitura-armairutxoak}
\textsuperscript{(17)}: han dena ordenan dago.

%%K t06-tit-4 sub
\VUsaut{5.0mm}
\VUpk{10.2pt}{40}{Bihurria.}
\VUsaut{3.0mm}

%%K t06-c2-07-1 p
7. --- Haren neba Paulok \VUgras{liburu} bat \textsuperscript{(19)} bilatzen du
\VUgras{liburu-armairuan} \textsuperscript{(18)}. \VUgras{Mutil}
\textsuperscript{(20)} ezin geldi egon daitekeena eta erabat jasanezina da.
Armairutik zintzilik, goregi dagoen liburu bat hartzeko, gainean dagoen
\VUgras{bustoa} \textsuperscript{(21)} eraisteko arriskua du.

%%K t06-c2-08-1 p
8. --- Zorakeria hori egiteko baliatzen da ama solasean dagoela bere adiskide
batekin, \VUgras{dama} batekin \textsuperscript{(22)}, zeina egun batzuk haren
etxean igarotzera etorri baita. Ama \VUgras{dibanean} \textsuperscript{(24)}
eserita dago, \VUgras{besaulkiaren} \textsuperscript{(25)} ondoan, eta haren
adiskidea \VUgras{aulki} batean \textsuperscript{(27)}, zeinetik
\VUgras{bizkarraldea} \textsuperscript{(26)} baino ez baitzaigu ikusten.

%%K t06-c2-09-1 p
9. --- Hormako \VUgras{marko} handian \textsuperscript{(30)} ahaide baten
\VUgras{erretratua} \textsuperscript{(29)} dago. Mahai gainean dagoen
\VUgras{albumean} \textsuperscript{(32)} gainerako senideen \VUgras{argazkiak}
\textsuperscript{(33)} bildu dira. Haurrak orain arte begiratzen aritu dira, oraindik
\VUgras{aulkiak} \textsuperscript{(31)} mahaiaren inguruan bilduta baitaude.

%%K t06-c2-10-1 p
10. --- Portzelanazko \VUgras{ontzi txinatarra} \textsuperscript{(34)},
\VUgras{paper-labanaren} \textsuperscript{(35)} ondoan, Mariari oparitu zioten
izen-egunean, eta gelaren izkina hartzen duen \VUgras{biombua}
\textsuperscript{(36)} haren osabak ekarri zuen Txinatik.

%%K t06-c2-11-1 p
11. --- \VUgras{Koadro} ederrez \textsuperscript{(37)} alaitua,
\VUgras{hormakoez} \textsuperscript{(42)}, sabaiko \VUgras{argi-armiarmaz}
\textsuperscript{(38)} argiztatua, zeinaren \VUgras{bonbilla elektrikoek}
\textsuperscript{(39)} hain alai argitzen baitute arratsetan, egongela hau oso
atsegina da. Zoruan zabaldutako \VUgras{tapiz} \textsuperscript{(40)} epelak eta
hain \VUgras{txirrin elektriko} erosoak \textsuperscript{(41)} erabat gozo
bihurtzen dute.

%%K t06-tit-5 sub
\VUsaut{3.6mm}
\VUpk{10.2pt}{40}{Jangela.}
\VUsaut{3.0mm}

%%K t06-c3-01-1 p
1. --- Bazkarira deitu dute. Familia osoa, Maria izan ezik, mahaiaren inguruan
bildu da. Jangela atsegin epelduta dago \VUgras{zeramikazko labe} bikainaz
\textsuperscript{(1)}, bere \VUgras{hodi} finarekin \textsuperscript{(2)}.

%%K t06-c3-02-1 p
2. --- Gela honetan altzari gutxi dago. Hormaren ondoan, \VUgras{taburetearen}
\textsuperscript{(4)} gainean, \VUgras{eskuak garbitzeko iturritxo} dotorea
\textsuperscript{(3)} dago zintzilik. Ondoan \VUgras{aparadorea}
\textsuperscript{(5)} dago, zeinaren gainean plater hotzak, postreak eta oraindik
behar ez den ontziteria jartzen baitira.

%%K t06-c3-03-1 p
3. --- Hala, goiko apalean \VUgras{oilasko errea} \textsuperscript{(7)},
\VUgras{arraina} \textsuperscript{(8)} eta \VUgras{ontzi} bat
\textsuperscript{(10)} \VUgras{frutaz} \textsuperscript{(9)} betea ikusten
ditugu.
Beherago \VUgras{gazta} \textsuperscript{(11)}, \VUgras{ogia}
\textsuperscript{(12)} eta \VUgras{ontzi-multzoa} \textsuperscript{(13)} daude,
zeinetan entsalada ontzeko behar den guztia baitago: olioa, ozpina,
\VUgras{gatza} \textsuperscript{(14)} eta \VUgras{piperbeltza}
\textsuperscript{(15)}. Azkenik, beheko apalean \VUgras{pastela}
\textsuperscript{(17)} dago \VUgras{mostaza-ontziaren} \textsuperscript{(16)}
ondoan.

%%K t06-c3-04-1 p
4. --- Jangelako erlojua, \VUgras{hormakoa} \textsuperscript{(20)} deitzen
dena, forma oso bitxikoa da. Zurezko marko batean sartuta dago, zeinetan
\VUgras{barometroa} \textsuperscript{(19)} eta \VUgras{termometroa}
\textsuperscript{(18)} ere baitaude.

%%K t06-tit-6 sub
\VUsaut{4.6mm}
\VUpk{10.2pt}{40}{Nola ematen den bazkaria.}
\VUsaut{3.0mm}

%%K t06-c3-05-1 p
5. --- Gela osoaren inguruan hormak \VUgras{zurezko panelez}
\textsuperscript{(21)} estalita daude, argizaritutako haritzezkoak. Oihalezko
\VUgras{errezel} batek \textsuperscript{(22)} nahiko ondo ixten du berandara
daraman atea. Hara joango da familia bazkalondoren kafea hartzera.
\VUgras{Katiluak} \textsuperscript{(24)} eta \VUgras{platertxoak}
\textsuperscript{(25)} \VUgras{beira-armairuan} \textsuperscript{(23)} jarrita
daude jada.

%%K t06-c3-06-1 p
6. --- Mahaiaren gainean sabaira \VUgras{zintzilikako lanpara} bat
\textsuperscript{(26)} lotuta dago; haren argi oso biziak jangela osoa betetzen du
arratsetan.

%%K t06-c3-07-1 p
7. --- Begira diezaiogun orain \VUgras{bazkal-mahaiari} berari
\textsuperscript{(27)}. Familia sartu zenean, mahaia jarrita zegoen.
\VUgras{Mahai-zapiaren} \textsuperscript{(28)} gainean leku bakoitzean
\VUgras{platera} \textsuperscript{(33)}, \VUgras{serbileta}
\textsuperscript{(29)}, \VUgras{koilara} \textsuperscript{(34)},
\VUgras{sardexka} \textsuperscript{(35)}, \VUgras{aiztoa} \textsuperscript{(36)}
eta \VUgras{edalontzia} \textsuperscript{(32)} jarri zituzten. Han zeuden
ardorako \VUgras{botilak} \textsuperscript{(30)} eta urerako \VUgras{karrafa}
\textsuperscript{(31)} ere.

%%K t06-c3-08-1 p
8. --- \VUgras{Zerbitzariak} \textsuperscript{(39)} lehenik \VUgras{zopa-ontzia}
\textsuperscript{(37)} ekarri du, zeinetan zopak lurruna botatzen baitu oraindik.
Orain lehen \VUgras{platera} \textsuperscript{(38)} ematen du, haragizkoa. Gero
denei eskainiko dizkie aipatu ditugunak; eta azkenik, \VUgras{postrearen}
ondoren, etxekoak berandara pasatu direla baliatuko du mahaia jasotzeko eta
jangela berriro txukuntzeko.

%%K t06-c3-08-2 p
\textit{Oharra.} --- \textit{Janariari buruzko eguneroko hitzak lerro
orokorretan bakarrik eman dira hemen. Zerrenda «Hotela» (13.) eta «Merkatua»
(15.) bi tauletan osatuko da.}

%%K t06-tit-7 sub
\VUsaut{4.6mm}
\VUpk{10.2pt}{40}{Sukaldea.}
\VUsaut{3.0mm}

%%K t06-c4-01-1 p
1. --- Sukaldean, ate erdi irekitik begiratzen dugun horretan, nahaste
pintoreskoa aurkitzen dugu. \VUgras{Sutegiaren} \textsuperscript{(1)} aurrean,
non \VUgras{sua} \textsuperscript{(3)} kraskatzen baita, \VUgras{burduntzi-makina}
\textsuperscript{(5)} jarri dute. \VUgras{Erre-haragi} bikaina
\textsuperscript{(7)} \VUgras{burduntzian} \textsuperscript{(6)} sartuta dago eta
egiten ari da.

%%K t06-c4-02-1 p
2. --- \VUgras{Hauspoa} \textsuperscript{(8)} eta \VUgras{ikatz-palatxoa}
\textsuperscript{(9)}, erabili berri dituztenak, zoruan geratu dira
\VUgras{enborrekin} \textsuperscript{(10)} batera.

%%K t06-c4-03-1 p
3. --- Tximiniako apalean ordena dago: \VUgras{farola} \textsuperscript{(11)},
\VUgras{pospolo-ontzia} \textsuperscript{(13)} bere \VUgras{pospoloekin}
\textsuperscript{(12)}, \VUgras{kandelargia} \textsuperscript{(14)} eta
\VUgras{gaueko kandelargia} \textsuperscript{(15)} bere \VUgras{kandelarekin}
\textsuperscript{(16)} --- dena bere lekuan. Aitzitik, \VUgras{ikatz-ontzi}
\textsuperscript{(18)} handia, \VUgras{ikatzez} \textsuperscript{(17)} betea,
zeinaren bila \VUgras{sukaldaria} \textsuperscript{(23)} sotora jaitsi berri
baita, sukaldearen erdian utzi dute.

%%K t06-c4-04-1 p
4. --- Kontua da emakume gaixoak presa duela, gosaria garaiz iritsi dadin. Orain
\VUgras{sukaldearen} \textsuperscript{(19)} ondoan dago, \VUgras{saltsa}
\textsuperscript{(24)} nahasten, zeinari ez baitzaio erretzen utzi behar.

%%K t06-c4-05-1 p
5. --- \VUgras{Labean} \textsuperscript{(20)} haurrentzat askaritarako egin duen
pastela erretzen ari da. Sukaldearen gainean \VUgras{burruntzalia}
\textsuperscript{(21)} eta \VUgras{bitsadera} \textsuperscript{(22)} daude
zintzilik.

%%K t06-tit-8 sub
\VUsaut{4.0mm}
\VUpk{10.2pt}{40}{Sukaldeko ontziteria.}
\VUsaut{3.0mm}

%%K t06-c4-06-1 p
6. --- \VUgras{Lapiko-multzoa} \textsuperscript{(25)}, gelaren barrenean
zintzilik dagoena, oso garbi dago. Kobrezko \VUgras{lapiko} ederrak
\textsuperscript{(26)}, \VUgras{esne-ontzia} \textsuperscript{(27)},
\VUgras{iragazkia} \textsuperscript{(28)} eta \VUgras{birringailua}
\textsuperscript{(29)} arretaz garbituta daude.

%%K t06-c4-07-1 p
7. --- \VUgras{Gatz-ontzia} \textsuperscript{(30)} apaltxo batean dago
\VUgras{te-ontziaren} \textsuperscript{(31)} eta \VUgras{olio-ontziaren}
\textsuperscript{(32)} ondoan. \VUgras{Kaxoian} \textsuperscript{(33)}, pentsatu
behar da, mahai-tresnak eta sukaldeko aiztoak daude.

%%K t06-c4-08-1 p
8. --- \VUgras{Erratza} \textsuperscript{(34)} eta \VUgras{lumazko astingailua}
\textsuperscript{(35)} izkinan daude. Sukaldariak \VUgras{ebakitzeko enborra}
\textsuperscript{(36)} erabili berri du; leihoaren ondoan utzi du.

%%K t06-c4-09-1 p
9. --- \VUgras{Eskuoihala} \textsuperscript{(37)} horman dago zintzilik
\VUgras{inbutuaren} \textsuperscript{(38)} ondoan. Sukaldearen zati honetan
\VUgras{parrila} \textsuperscript{(39)}, \VUgras{txikitzekoa}
\textsuperscript{(40)} eta \VUgras{ebakitzeko oholtxoa} \textsuperscript{(41)}
gordetzen dira. Gero, txikitzekoaren gainean, \VUgras{apalean}
\textsuperscript{(42)}, \VUgras{eltzeak} \textsuperscript{(43)},
\VUgras{burdinurtuzko lapikoa} \textsuperscript{(44)} bere \VUgras{estalkiarekin}
\textsuperscript{(45)} eta \VUgras{pertza} \textsuperscript{(46)} daude.
\VUgras{Zartagina} \textsuperscript{(47)} ondoan dago zintzilik.

%%K t06-c4-10-1 p
10. --- Kobrezko \VUgras{txorrotak} \textsuperscript{(48)} bakarrik ematen dio
distira izkina honi. \VUgras{Harraska} \textsuperscript{(49)}, zeinaren gainean
\VUgras{pitxer} handi bat \textsuperscript{(50)} baitago, oso erosoa izango da
bere armairutxoarekin, non \VUgras{ozpin-kupeltxoa} \textsuperscript{(51)} bere
\VUgras{txorrotatxoarekin} \textsuperscript{(52)}, \VUgras{zabor-ontzia}
\textsuperscript{(53)} eta \VUgras{ontzi zikinak} \textsuperscript{(54)}
gordetzen baitira.

%%K t06-tit-9 sub
\VUsaut{4.0mm}
\VUpk{10.2pt}{40}{Sukaldean gordetzen den beste zerbait.}
\VUsaut{3.0mm}

%%K t06-c4-11-1 p
11. --- \VUgras{Aska} handia \textsuperscript{(55)}, zeinetan \VUgras{barazkiak}
\textsuperscript{(58)} garbitzen baitira, eta burdinurtuzko \VUgras{pertza}
\textsuperscript{(56)}, \VUgras{zeramikazko zoruaren} \textsuperscript{(57)}
gainean dagoena, armairutxo horrenak dira ere, seguru asko.

%%K t06-c4-12-1 p
12. --- Sukaldariari sutegirik ez zaio falta, hementxe baitago, mahaiaren
izkinan, \VUgras{gas-erregailua} \textsuperscript{(59)}, eta haren gainean ura
berotzen ari da kazola batean. \VUgras{Lurruna} \textsuperscript{(60)}, handik
irteten dena, esaten digu ura irakiten egongo dela. Ur hori kaferako izango da,
seguru asko, mahaiaren beste muturrean \VUgras{kafe-ontzia}
\textsuperscript{(64)} eta \VUgras{kafe-errota} \textsuperscript{(63)} baitaude.

%%K t06-c4-13-1 p
13. --- Erregailua \VUgras{gas-adarrari} \textsuperscript{(62)} lotuta dago
\VUgras{goma-hodi} luze batez \textsuperscript{(61)}.

%%K t06-c4-14-1 p
14. --- \VUgras{Egun-langileak} \textsuperscript{(66)}, sukaldariari laguntzera
etortzen denak, barazkiak prestatzen ditu bazkarirako. Garbitu eta zuritu
dituenean, \VUgras{lapikoan} \textsuperscript{(65)} sartuko ditu, egosteko garaia
iritsi arte.

%%K t06-tit-10 sub
\VUsaut{4.0mm}
{\centering\fontsize{10.2pt}{10.2pt}\selectfont\textls[40]{\scshape Bainugela.} \textit{(Bainatzeko gela.)}\par}
\VUsaut{3.0mm}

%%K t06-c5-01-1 p
1. --- Leku bat besterik ez zaigu ikustea geratzen etxeko gela nagusiak
ezagutzeko: bainugelaren antolamendua. Ez da luzea izango, txikia baita, eta
laster ikusiko dugu \VUgras{bainuontziak} \textsuperscript{(1)}
\VUgras{berogailu} on bat duela \textsuperscript{(2)}, \VUgras{xaboi-ontzia}
\textsuperscript{(3)} bainatzen denaren eskura dagoela, eta \VUgras{eskuoihal-esekitokiak}
\textsuperscript{(5)} \VUgras{eskuoihalak} \textsuperscript{(4)} lehortzen
erraztu behar duela.

%%K t06-c5-02-1 p
2. --- Gela honen hormak \VUgras{beirazko lauzaz} \textsuperscript{(6)} estalita
daude, eta horri esker \VUgras{dutxa} \textsuperscript{(7)} jarri ahal izan dute,
erabiltzean zipriztintzen den urak hondatuko dituen beldurrik gabe. Zinkezko
\VUgras{ontzi} txiki bat \textsuperscript{(8)}, oin-bainuetarako, osatzen du
altzaritza.
\VUsaut{6.0mm}
\VUfilet{22mm}
